\section{Limitations}
\label{sec:limitations}

Our method equalizes local distributional solver defect rather than final
perceptual quality. It should therefore be interpreted as
teacher-alignment calibration, not as a direct optimizer for FID, CLIPScore,
ImageReward, or human preference. The matched SD1.5 and SDXL CFG sweeps are
retained because they show this distinction concretely.

The minimax result is asymptotic and relies on local power-law behavior. Very
small step counts such as $N=4$ or $N=8$ may fall outside the regime where the
theory is tight. The theory also assumes a deterministic ODE solver with
oracle-start defects; multistep and scheduler-history adapters require an
augmented-state analysis before the theorem applies exactly.

Calibration adds compute before generation. A schedule is practically useful
only when reuse amortizes the teacher and probe cost, so the experiments report
model-evaluation-equivalent calibration cost separately from generation cost.
Finally, the empirical monitor depends on the representativeness of the
calibration trajectories, the teacher solver accuracy, the chosen metric
$G(u)$, and the stability of the fitted order $q$.
